%======================================================================%
\section{Emergent Spacetime and Time}
\label{sec:spacetime}
%======================================================================%

Four-dimensional Lorentzian geometry is not a consequence of the $\E{6}$
coset metric alone. Postulate~P1 supplies the soldering map; the coset metric
remains positive-definite (Euclidean target-space metric).

%----------------------------------------------------------------------%
\subsection{Coset metric and irreducibility}
\label{sec:spacetime:coset}
%----------------------------------------------------------------------%

\begin{theorem}[Real-irreducibility of $\mathfrak{m}$]\label{thm:spacetime:irrep}
As a real $H$-module, $\mathfrak{m}=\mathbf{16}_{-3}\oplus\overline{\mathbf{16}}_{+3}$
is irreducible of \emph{complex type}: $\mathrm{End}_H(\mathfrak{m})\cong\C$.
\end{theorem}

\begin{proof}
Complexifying, $\mathfrak{m}_\C=\mathbf{16}_{-3}\oplus\overline{\mathbf{16}}_{+3}$
with non-isomorphic summands exchanged by conjugation. A real submodule would
be a complex submodule invariant under conjugation; Schur's lemma on each
summand forbids nontrivial proper submodules.
\end{proof}

\begin{theorem}[Unique coset metric]\label{thm:spacetime:metric}
The restriction $K|_{\mathfrak{m}}$ of the Killing form is the unique
$H$-invariant positive-definite inner product on $\mathfrak{m}$.
\end{theorem}

\begin{proof}
$H$-invariant bilinear forms correspond to $H$-invariant maps
$\mathfrak{m}\otimes\mathfrak{m}\to\R$. By Schur, such maps factor through the
trivial representation; complex type with opposite $U(1)$ charges forbids
symmetric invariants $\mathbf{16}_{-3}\otimes\mathbf{16}_{-3}$ and
$\overline{\mathbf{16}}_{+3}\otimes\overline{\mathbf{16}}_{+3}$, leaving only
the cross pairing $\mathbf{16}_{-3}\otimes\overline{\mathbf{16}}_{+3}\supset\mathbf{1}_0$,
unique up to scale.
\end{proof}

\begin{remark}
$K|_{\mathfrak{m}}$ is \textbf{positive-definite}, not Lorentzian. Signature
$(-,+,+,+)$ is carried by the soldering form, not by the Killing restriction.
\end{remark}

%----------------------------------------------------------------------%
\subsection{Postulate P1: soldering}
\label{sec:spacetime:soldering}
%----------------------------------------------------------------------%

The $2\times2$ Hermitian-octonionic block embeds as $H_2(\Oct)\cong\R^{1,9}$
with $\det$ the Minkowski norm; $\mathrm{SL}(2,\Oct)\cong\Spin(1,9)$.

\begin{postulate}[Soldering and dimension, P1]\label{post:spacetime:P1}
There exists a soldering map $\sigma$ identifying four clock-field directions
$e^a_\mu$ ($a=0,1,2,3$) with a four-dimensional subspace of $H_2(\Oct)$ carrying
$\eta_{ab}=\mathrm{diag}(-1,+1,+1,+1)$. The world manifold has $d=4$. The
vierbein emerges from the Maurer--Cartan $\mathfrak{m}$-component
(Section~\ref{sec:gauge-gravity}).
\end{postulate}

\begin{proposition}[Emergent vierbein]\label{prop:spacetime:vierbein}
For a coset section $g(x)\in\E{6}$ with $g^{-1}\partial_\mu g=A_\mu+e_\mu$,
the $\mathfrak{m}$-component defines a composite vierbein. Under P1,
$\eta_{\mu\nu}=e^a_\mu e^b_\nu\eta_{ab}$ has signature $(-,+,+,+)$.
\end{proposition}

%----------------------------------------------------------------------%
\subsection{Emergent time}
\label{sec:spacetime:time}
%----------------------------------------------------------------------%

\begin{proposition}[Page--Wootters clock]\label{prop:spacetime:time}
In the semiclassical regime of induced gravity
(Theorem~\ref{thm:gauge-gravity:planck}), the Hamiltonian constraint
$H_{\mathrm{phys}}|\psi\rangle=0$ holds on physical states. The soldered timelike
clock $\orderparam^0$ defines internal time; conditional evolution
$U(\tau)=e^{-iH_\tau\tau}$ on physical states recovers standard quantum dynamics.
\end{proposition}

\paragraph{Euclidean correlators and OS time.}
The coset metric $K|_{\mathfrak{m}}$ is positive-definite
(Theorem~\ref{thm:spacetime:metric}), so the kinetic sector of
\eqref{eq:dynamics:action} defines a Euclidean functional integral when the
world metric is taken Euclidean and the clock field $\orderparam^0$ is the
Osterwalder--Schrader (OS) time $\tau$. Lorentzian signature enters through
the soldering form (Postulate~\ref{post:spacetime:P1}), \emph{not} by Wick
rotating the coset metric.

\begin{definition}[Reflection positivity along $\orderparam^0$]
\label{def:spacetime:reflpos}
Let $\theta$ reflect $\tau\mapsto-\tau$ about a $\orderparam^0$-slice and
extend multiplicatively to field configurations. The Schwinger functions
$S_n$ satisfy \emph{reflection positivity} if, for finite supports
$F_i\subset\{\tau>0\}$,
\begin{equation}\label{eq:spacetime:reflpos}
  \sum_{i,j}\bar c_i\,c_j\;
  S\bigl(\theta F_i\cdot F_j\bigr)\;\ge\;0.
\end{equation}
\end{definition}

\begin{theorem}[OS reconstruction $\Rightarrow$ Lorentzian Wightman theory]
\label{thm:spacetime:OS}
Assume the Euclidean measure for \eqref{eq:dynamics:action} yields
Schwinger functions obeying Euclidean invariance on the $\tau$-slice,
reflection positivity \eqref{eq:spacetime:reflpos}, symmetry, and
clustering. Then the Osterwalder--Schrader reconstruction theorem
\cite{OsterwalderSchrader1973,OsterwalderSchrader1975} produces a Hilbert
space $\mathcal{H}$, a positive self-adjoint Hamiltonian $\widehat{H}\ge0$
generating $e^{-\widehat H\tau}$, and --- by analytic continuation
$\tau\to it$ in the Schwinger distributions --- Lorentzian Wightman
functions on $(\mathcal{M}^4,\eta^{(1,3)})$ with unitary evolution
$e^{-i\widehat H t}$ and reconstructed $\mathrm{SO}(1,3)$ covariance.
\end{theorem}

\begin{proposition}[Gaussian sector]
\label{prop:spacetime:rp-gaussian}
On a Euclidean background, the free $\sigma$-model with metric $K$ and
mass matrix $\Hess V(\orderparam_\star)$ satisfies
\eqref{eq:spacetime:reflpos}. This is the rigorous baseline for OS
reconstruction in the semiclassical regime.
\end{proposition}

\begin{postulate}[Interacting reflection positivity, P5]
\label{post:spacetime:P5}
The full interacting SJET $\sigma$-model, including the composite
$\Spin(10)\times U(1)$ gauge sector of
Section~\ref{sec:gauge-gravity}, satisfies reflection positivity
\eqref{eq:spacetime:reflpos} along the soldered clock $\orderparam^0$.
\end{postulate}

\begin{proposition}[No-go: RP is not derivable from P1 alone]
\label{prop:spacetime:rp-nogo}
Postulate~\ref{post:spacetime:P1} supplies a Lorentzian soldering form and
a Page--Wootters clock (Proposition~\ref{prop:spacetime:time}), but it does
\emph{not} determine the reflection positivity of the interacting Euclidean
measure. Without P5, or an equivalent lattice/nontrivial-field-theory proof,
unitary Lorentzian dynamics is conditional on
Theorem~\ref{thm:spacetime:OS}, not a consequence of the void--mirror axioms.
\end{proposition}

\begin{remark}[Status]
Theorem~\ref{thm:spacetime:OS} is rigorous given its hypotheses.
Proposition~\ref{prop:spacetime:rp-gaussian} covers the quadratic sector.
The gap to full unitary dynamics is precisely
Postulate~\ref{post:spacetime:P5}; see Open Problem~\textbf{O6} in
Section~\ref{sec:conclusions}.
\end{remark}